%%%%%%%%%%%%%%%%%%%%%%%%%%%%%%%%%%%%%%%%%
% LaTeX ашигласан тайлангийн загвар
%%%%%%%%%%%%%%%%%%%%%%%%%%%%%%%%%%%%%%%%%

%----------------------------------------------------------------------------------------
%	Хэрэглэх сангуудаа зарлах болон тохиргоонууд хийх хэсэг
%----------------------------------------------------------------------------------------

\documentclass[11pt, a4paper, oneside]{Thesis} % Цаасны хэмжээ, фонтын хэмжээ болон цаасныхаа хэдэн талд хэвлэхээ зааж өгнө

\graphicspath{{./Pictures/}} % Зургийн файлуудаа хаана хадгалахыг зааж өгнө

\usepackage[square, numbers, comma, sort&compress]
{natbib} % Ишлэл хийхдээ natbib пакежийг ашиглана 
\hypersetup{urlcolor=blue, colorlinks=true} % Вэб сайт рүү үсрэх холбоосны өнгийг өөрчилж өгнө
%\title{\ttitle} % Тайлангийн гарчигийг зааж өгнө
%\nocite{*} 
%----------------------------------------------------------------------------------------
%   Өөрийн мэдээллээ оруулах хэсэг
%----------------------------------------------------------------------------------------

\newcommand{\universityname}{монгол улсын их сургууль}
\newcommand{\facultyname}{Хэрэглээний Шинжлэх Ухаан, Инженерчлэлийн Сургууль}
\newcommand{\departmentname}{Электроник холбоо, инженерчлэлийн тэнхим}
\newcommand{\titlename}{Судалгааны ажлын нэр} 
\newcommand{\authorsname}{Овог нэр}
\newcommand{\supervisornameI}{Эрдмийн зэрэг, Овог нэр}
\newcommand{\supervisornameII}{Эрдмийн зэрэг, Овог нэр}
\newcommand{\cityname}{Улаанбаатар хот}
\newcommand{\thyear}{20..}
\newcommand{\CommChair}{Эрдмийн зэрэг, Овог нэр}
\newcommand{\CommClerk}{Эрдмийн зэрэг, Овог нэр}
\newcommand{\CommViceChair}{Эрдмийн зэрэг, Овог нэр}
\newcommand{\CommMemberI}{Эрдмийн зэрэг, Овог нэр}
\newcommand{\CommMemberII}{Эрдмийн зэрэг, Овог нэр}
\newcommand{\CommMemberIII}{Эрдмийн зэрэг, Овог нэр}
%\newcommand{\CommitteeMemberIV}{Комиссын гишүүний цол, нэр}
%\newcommand{\CommitteeMemberV}{Комиссын гишүүний цол, нэр}

\begin{document}

\frontmatter % Тайлангийн эхэн хэсгийн хуудасны дугаарлалтыг ромбо тоогоор хийнэ

\setstretch{1.3} % Мөр хоорондын зай 1.3

% Хуудасны толгой хэсгийн текстийг FancyHdr пакежийг ашиглан үүсгэнэ
\fancyhead{} % Бүх хөл болон толгой хэсгийн текстийг устгаж цэвэрлэнэ
\rhead{\thepage} % Хуудасны толгой хэсгийн текстийн баруун захад хуудасны дугаарыг байрлуулна
\lhead{} % Бүх хуудасны толгой хэсгийн зүүн захыг цэвэрлэнэ
\nocite{*}
\pagestyle{fancy} % Хуудасны толгой хэсгийн текстийг FancyHdr пакежийн загварын дагуу үүсгэнэ

% PDF файлын өгөгдлүүд
\hypersetup{pdftitle={\titlename}}
\hypersetup{pdfauthor=\authorsname}


%\input{./Chapters/Titlepage}
\input{./Chapters/DCCoverpage}
\input{./Chapters/DCTitlepage}
\input{./Chapters/Abstract}

%----------------------------------------------------------------------------------------
%	Гарчиг\Зургийн жагсаалт\Хүснэгтийн жагсаалтын хэсэг
%----------------------------------------------------------------------------------------

\pagestyle{fancy} % Энэ хэсгээс эхлэн хуудасны толгой хэсэгт текст байрлуулж эхлэнэ

\lhead{\emph{Гарчиг}} % Хуудасны толгой хэсгийн зүүн талын текстийг "Гарчиг" болгоно
\tableofcontents % Гарчигийн хуудсыг үүсгэнэ

\lhead{\emph{Зургийн жагсаалт}} % Хуудасны толгой хэсгийн зүүн талын текстийг "Зургийн жагсаалт" болгоно
\listoffigures % Зургийн жагсаалтын хэсгийг үүсгэнэ

\lhead{\emph{Хүснэгтийн жагсаалт}} % Хуудасны толгой хэсгийн зүүн талын текстийг "Хүснэгтийн жагсаалт" болгоно
\listoftables % Хүснэгтийн жагсаалтын хэсгийн үүсгэнэ

%----------------------------------------------------------------------------------------
%	Товчилсон үгсийн жагсаалт
%----------------------------------------------------------------------------------------

\clearpage % Шинэ хуудас үүсгэнэ

\setstretch{1.5} % Мөр хоорондын зайг 1.5 болгоно

\lhead{\emph{Товчилсон үгсийн жагсаалт}} % Хуудасны толгой хэсгийн зүүн талын текстийг "Товчилсон үгсийн жагсаалт" болгоно
\listofsymbols{ll} % 2 баганатай хүснэгтэй үүсгэнэ
{
\textbf{МУИС} & \textbf{М}онгол \textbf{У}лсын \textbf{И}х \textbf{С}ургууль \\
}


%----------------------------------------------------------------------------------------
%	Бүлгүд
%----------------------------------------------------------------------------------------

\mainmatter % Хуудасны дугаарлалтыг араб тоо болгоно(1,2,3)

\pagestyle{fancy} % Хуудасны толгой хэсгийн текстийг fancy хэлбэртэй болгоно.

% Бүлэг тус бүрийнхээ .tex файлыг Chapters хавтсанд хадгалж өгнө
% Дараах командаар бүлэг тус бүрийг оруулж өгнө
\input{./Chapters/Chapter1}
\input{./Chapters/Chapter2} 
\input{./Chapters/Chapter3}
\input{./Chapters/Chapter4} 
\input{./Chapters/Chapter5} 
%\input{./Chapters/Chapter6} 
%\input{./Chapters/Chapter7} 

%----------------------------------------------------------------------------------------
%	Хавсралт
%----------------------------------------------------------------------------------------

\addtocontents{toc}{\vspace{2em}} % Гарчигны хуудсанд зай үүсгэж өгнө

\appendix % Дараагийн бүлгүүд нь Хавсралт байх болно гэдгийг LaTeX-д хэлж өгнө

% Хавсралтуудаа тус бүрт нь .tex файл үүсгэн Appendices хавтсанд хадгалж өгөх хэрэгтэй.
% Appendices хавтсанд байгаа файлуудаа дараах командаар дуудаж өгнө.

\input{./Appendices/AppendixA}
%\input{./Appendices/AppendixB}
%\input{./Appendices/AppendixC}

\addtocontents{toc}{\vspace{2em}} % Гарчигны хуудсанд зай үүсгэж өгнө

\backmatter

%----------------------------------------------------------------------------------------
%	Ном зүй
\singlespacing
\bibliographystyle{IEEEtran}
\bibliography{Bibliography}

\end{document}  